% Source: https://artofproblemsolving.com/community/c3h2587629_8x8__13x5
% Stable identity: aops-topic-2587629-q01
Put the $5\times13$ rectangle at $0\le x\le13$, $0\le y\le5$. In the claimed arrangement, a triangle sits in the bottom left with vertices $(0,0)$, $(8,0)$, $(8,3)$, and a trapezoid sits in the bottom right with vertices $(8,0)$, $(13,0)$, $(13,5)$, $(8,3)$. The other two pieces are their images under the half-turn about the centre $\left(\dfrac{13}2,\dfrac52\right)$: a triangle $(13,5)$, $(5,5)$, $(5,2)$ and a trapezoid $(0,0)$, $(5,2)$, $(5,5)$, $(0,5)$.

The pieces would fill the rectangle only if their slanted edges formed the diagonal from $(0,0)$ to $(13,5)$. But the triangle's hypotenuse has slope $\dfrac38$, the trapezoid's slanted side has slope $\dfrac{5-3}{5}=\dfrac25$, and
\[\dfrac38<\dfrac5{13}<\dfrac25\qquad(0.375<0.3846<0.4).\]
So the lower pieces' top edge $(0,0)\to(8,3)\to(13,5)$ bends below the diagonal, and the upper pieces' bottom edge $(0,0)\to(5,2)\to(13,5)$ bends above it. Between them is a thin parallelogram with vertices $(0,0)$, $(8,3)$, $(13,5)$, $(5,2)$ that no piece covers. Its area, from the sides $(8,3)$ and $(5,2)$ at the origin, is
\[|8\cdot2-3\cdot5|=\boxed{1}.\]
That is exactly the missing unit: $64+1=65$. (The numbers $3,5,8,13$ are consecutive Fibonacci numbers, and $5\cdot13-8^2=1$ is an instance of Cassini's identity, which is why the error is so small.)

The cut square:

\begin{asy}
usepackage("amsmath");
size(3.5cm);
defaultpen(fontsize(10pt));
draw((0,0)--(8,0)--(8,8)--(0,8)--cycle,linewidth(0.8));
draw((0,3)--(8,3));draw((0,3)--(8,0));draw((5,3)--(3,8));
label("$8$",(4,0),S);label("$3$",(8,1.5),E);label("$5$",(8,5.5),E);
label("$5$",(2.5,3),N,fontsize(8pt));label("$3$",(6.5,3),N,fontsize(8pt));label("$3$",(1.5,8),S,fontsize(8pt));label("$5$",(5.5,8),S,fontsize(8pt));
\end{asy}

The rearranged pieces, to scale, with the uncovered sliver shaded:

\begin{asy}
usepackage("amsmath");
size(7cm);
defaultpen(fontsize(10pt));
pair O=(0,0),P=(8,3),Q=(13,5),R=(5,2);
fill(O--P--Q--R--cycle,gray);
draw((0,0)--(13,0)--(13,5)--(0,5)--cycle,linewidth(0.8));
draw(O--P--Q);draw(O--R--Q);draw(P--(8,0));draw(R--(5,5));
label("$8$",(4,0),S);label("$5$",(10.5,0),S);label("$3$",(8,1.5),W,fontsize(8pt));label("$3$",(5,3.6),E,fontsize(8pt));label("$5$",(13,2.5),E);
label("$(0,0)$",O,SW,fontsize(8pt));label("$(13,5)$",Q,NE,fontsize(8pt));
\end{asy}

The sliver stretched $25$ times across the diagonal (not to scale):

\begin{asy}
usepackage("amsmath");
size(6cm);
defaultpen(fontsize(10pt));
pair u=unit((13,5)),n=unit((-5,13));
pair T(pair p){return (dot(p,u),25*dot(p,n));}
pair O=(0,0),P=(8,3),Q=(13,5),R=(5,2);
fill(T(O)--T(P)--T(Q)--T(R)--cycle,lightgray);
draw(T(O)--T(P)--T(Q),linewidth(0.8));draw(T(O)--T(R)--T(Q),linewidth(0.8));draw(T(O)--T(Q),dotted);
dot(T(O));dot(T(P));dot(T(Q));dot(T(R));
label("$(0,0)$",T(O),W,fontsize(8pt));label("$(13,5)$",T(Q),E,fontsize(8pt));
label("$(8,3)$",T(P),S,fontsize(8pt));label("$(5,2)$",T(R),NE,fontsize(8pt));
label("slope $\dfrac38$",(T(O)+T(P))/2,S,fontsize(8pt));label("slope $\dfrac25$",(T(O)+T(R))/2,NW,fontsize(8pt));
\end{asy}
